% GENERATED FILE. Edit canonical lessons, then run npm run generate:books.
% canonical-source-hash: fnv1a64:b1a44c5c7b3a00a9
% canonical-lessons: KA-C11-bannagalu, KA-S120-letter-ha, KA-S142-digit-three

\chapter{Four Core Colours}
\label{ch:ka-core-colours}
\hlchaptermodality{11}

\begin{chapteropening}
\textbf{By the end of this chapter:} \emph{I can name four basic colours in Kannada — black, white, red, and blue.}

Say \kn{ಕಪ್ಪು}, \kn{ಬಿಳಿ}, \kn{ಕೆಂಪು} — the three native colours — then the borrowed \kn{ನೀಲಿ}, and point out that the week is Sanskritic where the colours are not.
\end{chapteropening}

\section[\texorpdfstring{\kn{ಕಪ್ಪು} \kn{ಬಿಳಿ} \kn{ಕೆಂಪು} \kn{ನೀಲಿ} (kappu biḷi kempu nīli)}{ಕಪ್ಪು ಬಿಳಿ ಕೆಂಪು ನೀಲಿ (kappu biḷi kempu nīli)}]{\kn{ಕಪ್ಪು}, \kn{ಬಿಳಿ}, \kn{ಕೆಂಪು}, \kn{ನೀಲಿ} — native colors, after a fully Sanskrit week}

\label{lesson:KA-C11-bannagalu}



\begin{quote}
Kannada's week ran entirely on Sanskrit deity-names. Its colors tell a different story — here, native Dravidian vocabulary is back in charge, except for one familiar exception.
\end{quote}

\subsection*{You'll want to know: Three native colors}
\begin{itemize}
\raggedright
  \item \textbf{\kn{ಕಪ್ಪು}} (\emph{kappu}) = \textquotedblleft{}\textbf{black}\textquotedblright{} — native Dravidian.
  \item \textbf{\kn{ಬಿಳಿ}} (\emph{biḷi}) = \textquotedblleft{}\textbf{white}\textquotedblright{} — native Dravidian.
  \item \textbf{\kn{ಕೆಂಪು}} (\emph{kempu}) = \textquotedblleft{}\textbf{red}\textquotedblright{} — native Dravidian.
\end{itemize}

\subsection*{You'll want to know: The familiar borrowed blue}
\textbf{\kn{ನೀಲಿ}} (\emph{nīli}) = \textquotedblleft{}\textbf{blue}\textquotedblright{} — from the \textbf{same Sanskrit} \kn{ನೀಲ} (\emph{nīla}) that Tamil, Malayalam, and Hindi all reach for. Even Kannada, whose week went fully Sanskritic, still keeps its everyday black/white/red \textbf{native} — it's specifically \textbf{blue} that pulls in the Sanskrit word, across all four Dravidian languages and Hindi alike.

\subsection*{Your turn}
\begin{itemize}
\raggedright
  \item \emph{Say it:} \textquotedblleft{}kappu, biḷi, kempu\textquotedblright{} — black, white, red
  \item \emph{Say it:} \textquotedblleft{}nīli\textquotedblright{} — blue, the shared Sanskrit word
  \item \emph{Say it:} the contrast — Kannada's week is Sanskritic, but its colors (except blue) are native
  \item \emph{Recall:} say \emph{ombattu}
  \item \emph{Read it:} \textbf{\kn{ಹತ್ತು}}
\end{itemize}

\begin{tcolorbox}[breakable,colback=teal!4,colframe=teal!35!black,title={Before you move on}]
Are Kannada's black/white/red native or Sanskrit? (\textbf{Native Dravidian} — kappu, biḷi, kempu.) Which color is the shared Sanskrit loan across Kannada, Tamil, Malayalam, and Hindi? (\textbf{Blue} — \kn{ನೀಲಿ} \emph{nīli} / \ta{நீலம்} \emph{nīlam} / \ml{നീല} \emph{nīla} / \dv{नीला} \emph{nīlā}.)
\end{tcolorbox}
\section[\texorpdfstring{\kn{ಹ} (ha)}{ಹ (ha)}]{\kn{ಹ} — one character, met inside words you already say}

\label{lesson:KA-S120-letter-ha}



\begin{quote}
Before the new one: \kn{◌ೋ} — what does it do?

One character this time. Just one — and you have been saying it for pages without knowing which mark on the page it was.
\end{quote}

\subsection*{Script you'll notice: \kn{ಹ}}
\textbf{\kn{ಹ}} — \emph{ha}.

It is a \textbf{consonant}, and in this script a consonant is never bare: it comes with an \emph{a} already in it. So it is not \emph{h}, it is \textbf{ha}.

You already say these, and every one of them has \kn{ಹ} somewhere inside it:

\begin{itemize}
\raggedright
  \item \textbf{\kn{ಹೌದು}} \emph{haudu} — yes (haudu)
  \item \textbf{\kn{ಹೆಸರು}} \emph{hesaru} — name
  \item \textbf{\kn{ಹೇಗೆ}} \emph{hēge} — how
  \item \textbf{\kn{ನೀವು} \kn{ಹೇಗಿದ್ದೀರಾ}?} \emph{nīvu hēgiddīrā} — how are you? (respectful)
\end{itemize}

\subsection*{Writing: \kn{ಹ}}
\hlblockfigure{\detokenize{figures/KA-S120-letter-ha-filmstrip.pdf}}{How \kn{ಹ} is written, stroke by stroke}

\begin{itemize}
\raggedright
  \item \textbf{1.} start at the upper right of the left ring and close it counterclockwise
  \item \textbf{2.} without lifting, arch over to the right and go down the right ring's outer side
  \item \textbf{3.} without lifting, round the right ring's base and climb its inner side to the waist
  \item \textbf{4.} without lifting, rise up the neck to the top bar
  \item \textbf{5.} lift, then draw the top bar from left to right
  \item \textbf{6.} without lifting, curl up into the hook
\end{itemize}

\textbf{Pen lifts: 1.} Both rings and the neck that rises between them are one run, starting at the top right of the left ring; the curled bar on top comes after the only lift.

\begin{quote}
This is one attested teaching order and not a national standard — Kannada handwriting is taught with school-to-school variation. Source: Gopala Krishna A, 'Kannada-alphabet-ha.gif', 40 frames, Wikimedia Commons, 25 May 2016, CC BY-SA 4.0.
\end{quote}

\subsection*{Your turn}
\begin{itemize}
\raggedright
  \item \emph{Look:} at these words, and find \kn{ಹ} in the ones that have it
\end{itemize}

\begin{quote}
\kn{ಹೌದು}  ·  \kn{ಹೆಸರು}  ·  \kn{ನಮಸ್ಕಾರ}
\end{quote}

\begin{itemize}
\raggedright
  \item \emph{Trace it:} \kn{ಹ} three times, saying \emph{ha} as you finish each one
  \item \emph{Look:} back at any page of this chapter and find \kn{ಹ} once more
\end{itemize}

\begin{tcolorbox}[breakable,colback=teal!4,colframe=teal!35!black,title={Before you move on}]
Which character is this — \kn{ಹ}? What sound does it carry? (\emph{\textbf{ha}}.) Name one word you already say that contains it.
\end{tcolorbox}
\section[\texorpdfstring{\kn{೩} (3)}{೩ (3)}]{\kn{೩} — the digit for mūru}

\label{lesson:KA-S142-digit-three}



\begin{quote}
Before the new shape: \textbf{\kn{೨}} — which number is it?

One digit this time, and you already know its name.
\end{quote}

\subsection*{Script you'll notice: \kn{೩}}
\textbf{\kn{೩}} — \textbf{3}, the digit for \textbf{\kn{ಮೂರು}} (\emph{mūru}), \textquotedblleft{}three\textquotedblright{}.

The digits behave exactly like the ones you grew up with: ten shapes, and place value doing the rest.

\subsection*{Writing: \kn{೩} — follow the numbered strip}
\hlblockfigure{\detokenize{figures/KA-S142-digit-three-filmstrip.pdf}}{How \kn{೩} is written, stroke by stroke}

Follow the numbered strip: start where movement 1 begins, and let each panel show you the next movement before your pen makes it. Then write \kn{೩} yourself — slowly, and larger than it is printed.

\begin{quote}
The strip shows \textbf{where to start the character and which way to travel}, in one attested order. That is taught with real variation from school to school, so the strip names its source beneath it: treat it as a sound way in, not the only one.
\end{quote}

\subsection*{Your turn}
\begin{itemize}
\raggedright
  \item \emph{Look:} at the digits you have met, and name each one aloud
\end{itemize}

\begin{quote}
\kn{೧}  ·  \kn{೨}  ·  \kn{೩}
\end{quote}

\begin{itemize}
\raggedright
  \item \emph{Say it:} \textquotedblleft{}mūru\textquotedblright{} as you point at \kn{೩}
\end{itemize}

\begin{tcolorbox}[breakable,colback=teal!4,colframe=teal!35!black,title={Before you move on}]
Which number is \textbf{\kn{೩}}? (\textbf{3} — \emph{mūru}, \kn{ಮೂರು}.) Name every digit you have met so far, in order.
\end{tcolorbox}
